\section{Conclusion}

Adding a finite speed of information to the self-replicating slingshot model turns a
convenient assumption into a measurable cost. Over a paired ensemble of identical galaxies,
light-speed lag slows the time to fill the field by a median of about 0 percent under powered
flight, about 30 percent under nearest-star slingshots, and about 51 percent under
maximum-boost slingshots; the slingshot penalties are positive in every one of the 32 seeds,
hold across coverage fractions and across a four-fold range of field sizes, and every case
still reaches full coverage. The dimensionless ratio $\Lambda = v/c$ gates whether the tax
appears at all, and at the powered cruise speed it does not; but because that ratio is the
same on every hop it does not predict the penalty's size. That is decided by how non-local
each hop is: fast-but-local flight is nearly immune, while the cost falls on the long-range
hops that made the swarm fast in the first place, whose wasted trips we measure to be about
twice as long. The slingshot speed-up survives contact with the finite speed of information,
but a real swarm pays a coordination tax that a perfect-information model hides, and pays it
in proportion to its reach. Two extensions would tighten the estimate without changing its
shape: knowledge gained in flight, and probe-to-probe relay of the settled map. Both would
lower the tax; neither can remove the physical floor beneath it.

\section*{Data Availability}
The model is a deterministic, seeded fold; every figure, table, and reported statistic in
this paper regenerates from source by running the swarm module's experiment and figure
scripts, and every quantitative claim traces to a cited source in the project's shared
bibliography. The code and its citation metadata are archived with the concept DOI
\texttt{10.5281/zenodo.21249296}.

\section*{Acknowledgment}
This work used no external funding.
